% ！TeX program = luatex
\documentclass[a4paper]{article}
% \documentclass[UTF8]{ctexart}
% 导言区
% biber bibtex
% \usepackage[backend=biber]{biblatex}

% \usepackage{CJKutf8} % 支持中
% fontset = none 消除Package ctex: Redefining CJKfamily `\CJKrmdefault'(ctex)	(FandolSong-Regular).
% 重新定义字体警告
\usepackage[fontset = none]{ctex} % 支持中文 必须先加载(LuaLaTeX)
% \setmonofont[Mapping={}]{Monaco}    %英文引号之类的正常显示，相当于设置英文字体
% \setsansfont{Monaco} %设置英文字体 Monaco, Consolas,  Fantasque Sans Mono
% \setmainfont{Monaco} %设置英文字体
% \setCJKmainfont{微软雅黑}  %中文字体设置
% \setCJKsansfont{楷体} %设置中文字体
% \setCJKmonofont{} %设置中文字体
% \usepackage{xeCJK} % 支持中文 for XeLaTeX
\usepackage{fontspec} % 支持中文 for LuaLaTeX
\setCJKmainfont{LXGW WenKai} % 设置中文字体
\setCJKsansfont{LXGW WenKai Mono} % 设置中文字体
\setCJKmonofont{LXGW WenKai Mono} % 设置中文字体
% 使用近似的替代字体
\newfontface\myitalic{LXGW WenKai Mono}[FakeSlant=0.25] % 模拟斜体

% 设置emoji字体 只支持 LuaLaTeX
% 其他编译器用Noto Color Emoji只会报错
% 只能用Noto Emoji
\newfontfamily{\emoji}{Noto Emoji}

\usepackage{babel} % load before hyperref
\usepackage[colorlinks, linkcolor=skyblue]{hyperref} % 支持超链接式

\babelprovide[import, main]{english} % 加载英语支持
\babelprovide[import]{chinese-simplified} % 加载简体中文并设为主语言



\usepackage{amsmath} % 支持数学符号 align*
\usepackage{amssymb} % 更多数学符号 包括 ∄ 等
\usepackage{esint} % 支持多重积分 包括 \unicode{222F} ∮∯ 等
\usepackage{unicode-math}
\unimathsetup{math-style=TeX} % 使用传统的TeX数学样

\usepackage{listing} % 支持代码高亮，设置高亮样式

% 处理对齐问题：\toprule， \midrule

%导言区插入下面三行
\usepackage{graphicx} %插入图片的宏包
\usepackage{float} %设置图片浮动位置的宏包
% Deprecated 不再更新
% \usepackage{subfigure} %插入多图时用子图显示的宏包
% 使用 subcaption 替代
\usepackage{subcaption} % 支持子图
% https://xfxuezhang.blog.csdn.net/article/details/144632165?spm=1001.2101.3001.6650.4&utm_medium=distribute.pc_relevant.none-task-blog-2%7Edefault%7EYuanLiJiHua%7EPosition-4-144632165-blog-103762494.235%5Ev43%5Epc_blog_bottom_relevance_base7&depth_1-utm_source=distribute.pc_relevant.none-task-blog-2%7Edefault%7EYuanLiJiHua%7EPosition-4-144632165-blog-103762494.235%5Ev43%5Epc_blog_bottom_relevance_base7&utm_relevant_index=9
% 还是使用 subfigure {begin/end}

\usepackage{pdfpages} % 支持插入PDF
% 和pdfpages里面的xcolor包冲突
% \usepackage{xcolor} % 支持颜色设置
\usepackage{tabularx} % 支持表格自动列宽、

\usepackage{listings}   % 支持代码高亮，设置高亮样式
% \usepackage[usenames,dvipsnames]{xcolor}    % 支持颜色设置，用于高亮不同颜色关键字
\usepackage{embrac}   % 支持斜体直体 粗体等字体样式的宏包
\usepackage{colortbl}  % 支持颜色

\usepackage{booktabs} % 表格的支持需要加载宏包{booktabs}

\usepackage{hhline} % 支持表格分割线双横线双竖线
\usepackage{makecell} % 支持表格内换行
\usepackage{multirow} % 支持跨行合并单元格
\usepackage{multicol}
\usepackage{stfloats} % 防止表格乱跑


\usepackage{longtable} % 跨页表格

\usepackage{diagbox} % 支持表格斜线

\usepackage{ulem} % 支持删除线 https://www.jianshu.com/p/a1838fa53882

\usepackage{pifont} % 支持带圈数字

\usepackage{enumitem} % 支持列表缩进

\usepackage{bigstrut}

\usepackage[fontsize=8pt]{fontsize}
\usepackage{geometry} % 支持调整页边距

\usepackage{epsfig} % 支持插入eps图片
\usepackage{svg} % 支持插入svg图片

\usepackage{bookmark} % Fixes the rerunfilecheck issue

\svgsetup{
  inkscapeexe = /usr/bin/inkscape,
}

% ---------------- RISV-V Assembler syntax and style for latex lstlisting package ----------------
% RISC-V Assembler syntax and style for latex lstlisting package
% 
% These are risc-v commands as per our university (University Augsburg, Germany) guidelines.
%
% Author: Anton Lydike
%
% This code is in the public domain and free of licensing

% language definition
% \begin{Language_Definition}
\lstdefinelanguage[RISC-V]{Assembler}
{
  alsoletter={.}, % allow dots in keywords
  alsodigit={0x}, % hex numbers are numbers too!
  morekeywords=[1]{ % instructions
    lb, lh, lw, lbu, lhu,
    sb, sh, sw,
    sll, slli, srl, srli, sra, srai,
    add, addi, sub, lui, auipc,
    xor, xori, or, ori, and, andi,
    slt, slti, sltu, sltiu,
    beq, bne, blt, bge, bltu, bgeu,
    j, jr, jal, jalr, ret,
    scall, break, nop
  },
  morekeywords=[2]{ % section*s of our code and other directives
    .align,.ascii,.asciiz,.byte,.data,.double,.extern,
    .float,.globl,.half,.kdata,.ktext,.set,.space,.text,.word
  },
  morekeywords=[3]{ % registers
    zero, ra, sp, gp, tp, s0, fp,
    t0, t1, t2, t3, t4, t5, t6,
    s1, s2, s3, s4, s5, s6, s7, s8, s9, s10, s11,
    a0, a1, a2, a3, a4, a5, a6, a7,
    ft0, ft1, ft2, ft3, ft4, ft5, ft6, ft7,
    fs0, fs1, fs2, fs3, fs4, fs5, fs6, fs7, fs8, fs9, fs10, fs11,
    fa0, fa1, fa2, fa3, fa4, fa5, fa6, fa7
  },
  morecomment=[l]{;},   % mark ; as line comment start
  morecomment=[l]{\#},  % as well as # (even though it is unconventional)
  morestring=[b]",      % mark " as string start/end
  morestring=[b]'       % also mark ' as string start/end
}
% \end{Language_Definition}

% \begin{Language_Style_Setting}
\definecolor{myGreen}{RGB}{50, 150, 50} % 自定义颜色代Block 
% -------------------listing 设置，代码设置-------------------
\lstset{ %
    backgroundcolor=\color{white},   % choose the background color; you must add \usepackage{color} or \usepackage{xcolor}
    basicstyle=\footnotesize,        % the size of the fonts that are used for the code
    breakatwhitespace=false,         % sets if automatic breaks should only happen at whitespace
    breaklines=true,                 % sets automatic line breaking
    captionpos=bl,                    % sets the caption-position to bottom
    commentstyle=\color{myGreen},    % comment style
    deletekeywords={\ldots},
    % keeps spaces in text, useful for keeping indentation of code (possibly needs columns=flexible)
    % keywordstyle=\color{blue},       % keyword style
    % language=Python,                      % the language of the code
    morekeywords={*,\ldots},            % if you want to add more keywords to the set
    numbers=left,                    % where to put the line-numbers; possible values are (none, left, right)
    numbersep=5pt,                   % how far the line-numbers are from the code
    numberstyle=\tiny\color{gray}, % the style that is used for the line-numbers
    rulecolor=\color{black},         % if not set, the frame-color may be changed on line-breaks within not-black text (e.g. comments (green here))
    showspaces=false,                % show spaces everywhere adding particular underscores; it overrides 'showstringspaces'
    showstringspaces=false,          % underline spaces within strings only
    showtabs=false,                  % show tabs within strings adding particular underscores
    stepnumber=1,                    % the step between two line-numbers. If it's 1, each line will be numbered
    stringstyle=\color{orange},     % string literal style
    tabsize=4,                       % sets default tabsize to 2 spaces
    language=[RISC-V]Assembler,                   % all code is RISC-V
    identifierstyle=\color{teal},                 % user declared addresses are teal
    %keywordstyle=[1]\color{blue!80!black},        % instructions are blue
    %keywordstyle=[2]\color{orange!80!black},      % section*s/other directives are orange
    %keywordstyle=[3]\color{red!50!black},         % registers are red
    keywordstyle=[1]\color{blue},        % instructions are blue
    keywordstyle=[2]\color{orange},      % section*s/other directives are orange
    keywordstyle=[3]\color{red}         % registers are red
    %title=myPython.py                   % show the filename of files included with \lstinputlisting; also try caption instead of title
}
% ----------------
% \end{Language_Style_Setting}

% \setCJKmonofont{Adobe Fangsong Std} % 设置等宽字体



\begin{document}
\begin{align*}
    Result &=
    \frac{
        1
    }{
    e ^ 
    {
    -k \cdot
    \left[
    \begin{aligned}
        &\quad F_{CenterSurroundWeightedContrast}
        \\
        &\pm F_{CenterWeightedContrast} 
    \end{aligned}
    \circ ImageOrAnySource
    \right]
    - b
    }
}
\cdot 255
\end{align*}
\begin{align*}
    Result &=
    \frac{
        1
    }{
    e ^ 
    {
        -k \cdot
        \left[
        \begin{aligned}
            &\Sigma_{j \in Surround[1,2,3,4,5,6,7,8]}
            \left(
                \begin{aligned}
                    &\left[
                        F_{pdf}\left(
                            \Vert{
                                \overrightarrow{Block_{i=Center}}
                                - \overrightarrow{Block_{j \in Surround}}
                            }\Vert_{NormDimension}
                        \right) \sim \mathcal{N}(\mu, \sigma^{2})
                    \right]
                    \\
                    &\cdot
                    \left[
                        F_{Contrast}(F_{Attribute}(Block_{i =Center})， F_{Attribute}(Block_{j \in Surround}))
                    \right]
                \end{aligned}
            \right)
            \\
            &\pm F_{CenterWeightedContrast} 
        \end{aligned}
        \circ ImageOrAnySource
        \right]
        - b
        }
    }
    \cdot 255
\end{align*}
\begin{align*}
F_{CentralWeightedContrast} = 
    \left(
    \begin{aligned}
    &
        \Sigma_{(k, l)\ in\ \{(1, 8), (2, 7), (3, 6), (4, 5)\}}
    \left(
        \begin{aligned}
            &\left\{
                \begin{aligned}
                    &\left[
                        F_{pdf}\left(
                            \Vert{
                                \overrightarrow{Block_{k}}
                                - \overrightarrow{Block_{l}}
                            }\Vert_{NormDimension}
                        \right) \sim \mathcal{N}(\mu, \sigma^{2})
                    \right]
                    \\
                    &\left[
                        F_{pdf}\left(
                            \Vert{
                                \frac{
                                    \overrightarrow{Block_{k}}
                                    - \overrightarrow{Block_{l}}
                                }{2}
                            }\Vert_{NormDimension}
                        \right) \sim \mathcal{N}(\mu, \sigma^{2})
                    \right] \cdot 2
                    \\
                    &\left[
                        F_{pdf}\left(
                            \Vert{
                                \frac{
                                    \overrightarrow{Block_{Center}}
                                    - \overrightarrow{Block_{Center}}
                                }{2}
                            }\Vert_{NormDimension}
                        \right) \sim \mathcal{N}(\mu, \sigma^{2})
                    \right] \cdot 2
                \end{aligned}
            \right.
        \end{aligned}
    \right) \cdot 2
    \\
    &\cdot
    \left[
        F_{Contrast}(F_{Attribute}(Block_{k})， F_{Attribute}(Block_{l}))
    \right]
    \end{aligned}
    \right)
\end{align*}
\begin{align*}
    F_{\ldots WeighedContrastFunction}
    &=
    [F_{Weight} \sim Distribution_{AnyDistribution}(\cdots)]
    \cdot [F_{Contrast}]
\end{align*}
Note:
\begin{align*}
    \left\{
        \begin{aligned}
            &Gradient = k
            \\
            &Contrast = c
            \\
            &F_{CenterWeightedContrast}
            \\
            &F_{CenterSurroundContrast}
            \\
            &NormDimension
            \\
            &Distribution_{AnyDistribution}
        \end{aligned}
    \right.
\end{align*}
Could be changed like buiding blocks.
\end{document}